\documentclass{article}\begin{document}At a fixed base object and for an arbitrary ordinary payload transformation, evaluationPayload gives an isomorphism of functors from actual weighted single coproduct-span evaluation to the actual tagged iterate value. The two paths use the equivalence functors of All.evaluation for the two payloads, the change-of-payload functor on the single weighted endpoint, and the tagged changes of the actual iterate payload. It is not an equality of inverse functors: inverseSquare uses the unit and counit of the chosen equivalences to turn the forward comparison of \#196 into a square for its inverse. The remaining tagged square is assembled on actual tagged arrows, using the identity-span payload law at zero and the positive evaluation law at successors. Their paste is the claimed square; no assumption on the mate or on individual comparison components is added. Separately taggedBase assembles the genuine chosen native base squares for every length, including the identity-span chosen pull at zero, into a natural isomorphism against the pre-existing iterate packedMap. Its source is still the tagged sum of weighted fibres, not the single span. To obtain a base square for All.evaluation, first provide the single coproduct source cleavage by tagging component lifts, then compare the induced weighted pulls across the actual compare functor by factorization through the selected cartesian lifts. Merely conjugating the tagged pull across compare would define a different action and does not identify an independently chosen single-span pull. Payload identity/composition coherence of these new isomorphisms, assembled base identity/composition, mixed coherence, universe widening, flattening and algebra remain outside this result. Common base/grade/payload universe restrictions for the all-length telescope remain exactly those of \#196; the generic tagged and inverse-square lemmas are independently universe-polymorphic. Consult the pinned manuscript 458ff36a, lines 929--953, together with the definition of All.evaluation for the order and orientation of the two component comparisons.\end{document}
